\begin{problem}[VI.32.P20: Legacy DG4 Problem 13 --- unions of three lines]\label{prob:vi32-20}
Assume $k$ is algebraically closed and, to match the earlier running example, $\operatorname{char}k\neq 2$.  Compare
\[
 X=V(xy,xz,yz)\subset\mathbb A^3_k
 \qquad\text{and}\qquad
 Y=V(uv(u-v))\subset\mathbb A^2_k.
\]
Both are unions of three affine lines meeting at one point.  Prove that $X$ and $Y$ are not isomorphic by comparing their Zariski tangent spaces at the common intersection point.

\textbf{Provenance.} Canonical recovery of legacy DG4 Problem 13, whose tangent-space invariant was reserved for VI/32.  The two line-union models are the running examples already established in VI/05.
\begin{solution}
Each scheme has three irreducible components, all meeting only at the origin, so the component-incidence picture alone does not distinguish them.  At the origin of $X$, every generator $xy,xz,yz$ has zero linear part; equivalently the Jacobian matrix is zero there.  Hence
\[
 T_0X\cong k^3,
 \qquad \dim_kT_0X=3.
\]
For $Y$, the defining polynomial $uv(u-v)$ also has zero linear part at the origin, but $Y$ lies in the plane, so
\[
 T_0Y\cong k^2,
 \qquad \dim_kT_0Y=2.
\]
The origin is the unique point at which three irreducible components meet in either scheme, so any isomorphism would have to send origin to origin.  An isomorphism induces an isomorphism of local rings and therefore of cotangent spaces $\mathfrak m/\mathfrak m^2$, hence also of their dual tangent spaces.  The tangent dimensions $3$ and $2$ differ, so no isomorphism can exist.
\end{solution}
\end{problem}
